% GENERATED FILE. Edit canonical lessons, then run npm run generate:books.
% canonical-source-hash: fnv1a64:c1cbac60ccef588c
% canonical-lessons: UR-C318-namaz, UR-C318-idi, UR-C318-sivayan, UR-C318-mubarak-bad, UR-C318-lipatna

\chapter{Eid Morning}
\label{ch:ur-eid-morning}
\hlchaptermodality{318}

\begin{chapteropening}
\textbf{By the end of this chapter:} \emph{I can tell how a family spent a festival morning.}

Complete the last lesson of chapter 318: I can tell how a family spent a festival morning.
\end{chapteropening}

\section[\texorpdfstring{\ur{نماز} (namāz)}{نماز (namāz)}]{\ur{نماز} (namāz) — the set prayer}

\label{lesson:UR-C318-namaz}



\begin{quote}
Before the new one: say the Urdu for to glitter, then the Urdu for festive bustle.
\end{quote}

\subsection*{You'll want to know: \ur{نماز}}
\textbf{\ur{نماز}} — \emph{namāz} — \textquotedblleft{}the set prayer\textquotedblright{}.

\begin{quote}
\ur{صبح} \ur{سب} \ur{نے} \ur{نماز} \ur{پڑھی۔}
\end{quote}

\begin{quote}
\emph{subah sab ne namāz paṛhī.}
\end{quote}

\begin{quote}
In the morning everyone said the prayer.
\end{quote}

Feminine.

\begin{cousinweb}
\textbf{\ur{نماز}} (\emph{namāz}) is usually traced to Persian \emph{namāz}, \textquotedblleft{}a bowing, a prayer\textquotedblright{}. It is said, in Urdu, by reading it: \textbf{\ur{نماز} \ur{پڑھی}}.
\end{cousinweb}

\begin{grammarlens}[title={What you've built}]
One more word for telling what happened.
\end{grammarlens}

\subsection*{Your turn}
\begin{itemize}
\raggedright
  \item \emph{Say it:} \emph{namāz}
  \item \emph{Say it:} \emph{namāz}, once more
  \item \emph{Say it:} \emph{subah sab ne namāz paṛhī.} — In the morning everyone said the prayer.
  \item \emph{Recall:} say the Urdu for to glitter, then the Urdu for festive bustle, then say \emph{namāz} again
\end{itemize}

\begin{tcolorbox}[breakable,colback=teal!4,colframe=teal!35!black,title={Before you move on}]
What does \ur{نماز} mean? (\textquotedblleft{}The set prayer\textquotedblright{}.) Say it once more.
\end{tcolorbox}
\section[\texorpdfstring{\ur{عیدی} (īdī)}{عیدی (īdī)}]{\ur{عیدی} (īdī) — money given at Eid}

\label{lesson:UR-C318-idi}



\begin{quote}
Before the new one: say the Urdu for festive bustle, then the Urdu for the set prayer.
\end{quote}

\subsection*{You'll want to know: \ur{عیدی}}
\textbf{\ur{عیدی}} — \emph{īdī} — \textquotedblleft{}money given at Eid\textquotedblright{}.

\begin{quote}
\ur{بچوں} \ur{کو} \ur{دادا} \ur{سے} \ur{عیدی} \ur{ملی۔}
\end{quote}

\begin{quote}
\emph{bachchoṅ ko dādā se īdī milī.}
\end{quote}

\begin{quote}
The children got Eid money from their grandfather.
\end{quote}

Feminine.

\begin{cousinweb}
\textbf{\ur{عیدی}} (\emph{īdī}) is built on \textbf{\ur{عید}}, usually traced, through Persian, to Arabic \emph{ʿīd}, with the ending \emph{-ī}: the money elders give children on the day.
\end{cousinweb}

\begin{grammarlens}[title={What you've built}]
One more word for telling what happened.
\end{grammarlens}

\subsection*{Your turn}
\begin{itemize}
\raggedright
  \item \emph{Say it:} \emph{īdī}
  \item \emph{Say it:} \emph{īdī}, once more
  \item \emph{Say it:} \emph{bachchoṅ ko dādā se īdī milī.} — The children got Eid money from their grandfather.
  \item \emph{Recall:} say the Urdu for festive bustle, then the Urdu for the set prayer, then say \emph{īdī} again
\end{itemize}

\begin{tcolorbox}[breakable,colback=teal!4,colframe=teal!35!black,title={Before you move on}]
What does \ur{عیدی} mean? (\textquotedblleft{}Money given at Eid\textquotedblright{}.) Say it once more.
\end{tcolorbox}
\section[\texorpdfstring{\ur{سویاں} (sivayāṅ)}{سویاں (sivayāṅ)}]{\ur{سویاں} (sivayāṅ) — sweet vermicelli}

\label{lesson:UR-C318-sivayan}



\begin{quote}
Before the new one: say the Urdu for the set prayer, then the Urdu for money given at Eid.
\end{quote}

\subsection*{You'll want to know: \ur{سویاں}}
\textbf{\ur{سویاں}} — \emph{sivayāṅ} — \textquotedblleft{}sweet vermicelli\textquotedblright{}.

\begin{quote}
\ur{امی} \ur{نے} \ur{سویاں} \ur{پکائیں۔}
\end{quote}

\begin{quote}
\emph{ammī ne sivayāṅ pakāīṅ.}
\end{quote}

\begin{quote}
Mum cooked sweet vermicelli.
\end{quote}

Feminine, used in the plural.

\begin{cousinweb}
\textbf{\ur{سویاں}} (\emph{sivayāṅ}) is inherited: CDIAL lists the Hindi cousins \emph{sewaī̃} and \emph{sẽwiyā̃}, \textquotedblleft{}vermicelli\textquotedblright{}, under Sanskrit \emph{śamitā} (CDIAL 12307). Urdu says it in the plural: \textbf{\ur{سویاں} \ur{پکائیں}}.
\end{cousinweb}

\begin{grammarlens}[title={What you've built}]
One more word for telling what happened.
\end{grammarlens}

\subsection*{Your turn}
\begin{itemize}
\raggedright
  \item \emph{Say it:} \emph{sivayāṅ}
  \item \emph{Say it:} \emph{sivayāṅ}, once more
  \item \emph{Say it:} \emph{ammī ne sivayāṅ pakāīṅ.} — Mum cooked sweet vermicelli.
  \item \emph{Recall:} say the Urdu for the set prayer, then the Urdu for money given at Eid, then say \emph{sivayāṅ} again
\end{itemize}

\begin{tcolorbox}[breakable,colback=teal!4,colframe=teal!35!black,title={Before you move on}]
What does \ur{سویاں} mean? (\textquotedblleft{}Sweet vermicelli\textquotedblright{}.) Say it once more.
\end{tcolorbox}
\section[\texorpdfstring{\ur{مبارک} \ur{باد} (mubārak bād)}{مبارک باد (mubārak bād)}]{\ur{مبارک} \ur{باد} (mubārak bād) — congratulations}

\label{lesson:UR-C318-mubarak-bad}



\begin{quote}
Before the new one: say the Urdu for money given at Eid, then the Urdu for sweet vermicelli.
\end{quote}

\subsection*{You'll want to know: \ur{مبارک} \ur{باد}}
\textbf{\ur{مبارک} \ur{باد}} — \emph{mubārak bād} — \textquotedblleft{}congratulations\textquotedblright{}.

\begin{quote}
\ur{پڑوسیوں} \ur{نے} \ur{ہمیں} \ur{عید} \ur{کی} \ur{مبارک} \ur{باد} \ur{دی۔}
\end{quote}

\begin{quote}
\emph{paṛosiyoṅ ne hameṅ īd kī mubārak bād dī.}
\end{quote}

\begin{quote}
The neighbours wished us a happy Eid.
\end{quote}

Feminine.

\begin{cousinweb}
\textbf{\ur{مبارک} \ur{باد}} (\emph{mubārak bād}) is usually traced to Arabic \emph{mubārak}, \textquotedblleft{}blessed\textquotedblright{}, and Persian \emph{bād}, \textquotedblleft{}may it be\textquotedblright{}, a form of the Persian verb \textquotedblleft{}to be\textquotedblright{}: \textquotedblleft{}may it be blessed\textquotedblright{}. It is feminine, and the one wished takes \textbf{\ur{کو}}: \textbf{\ur{ہمیں} \ur{مبارک} \ur{باد} \ur{دی}}.
\end{cousinweb}

\begin{grammarlens}[title={What you've built}]
One more word for telling what happened.
\end{grammarlens}

\subsection*{Your turn}
\begin{itemize}
\raggedright
  \item \emph{Say it:} \emph{mubārak bād}
  \item \emph{Say it:} \emph{mubārak bād}, once more
  \item \emph{Say it:} \emph{paṛosiyoṅ ne hameṅ īd kī mubārak bād dī.} — The neighbours wished us a happy Eid.
  \item \emph{Recall:} say the Urdu for money given at Eid, then the Urdu for sweet vermicelli, then say \emph{mubārak bād} again
\end{itemize}

\begin{tcolorbox}[breakable,colback=teal!4,colframe=teal!35!black,title={Before you move on}]
What does \ur{مبارک} \ur{باد} mean? (\textquotedblleft{}Congratulations\textquotedblright{}.) Say it once more.
\end{tcolorbox}
\section[\texorpdfstring{\ur{لپٹنا} (lipaṭnā)}{لپٹنا (lipaṭnā)}]{\ur{لپٹنا} (lipaṭnā) — to cling; to hug}

\label{lesson:UR-C318-lipatna}



\begin{quote}
Before the new one: say the Urdu for sweet vermicelli, then the Urdu for congratulations.
\end{quote}

\subsection*{You'll want to know: \ur{لپٹنا}}
\textbf{\ur{لپٹنا}} — \emph{lipaṭnā} — \textquotedblleft{}to cling; to hug\textquotedblright{}.

\begin{quote}
\ur{بچی} \ur{ماں} \ur{سے} \ur{لپٹ} \ur{گئی۔}
\end{quote}

\begin{quote}
\emph{bachchī māṅ se lipaṭ gaī.}
\end{quote}

\begin{quote}
The little girl clung to her mother.
\end{quote}

Past (she …): \ur{وہ} \ur{لپٹ} \ur{گئی}, \emph{voh lipaṭ gaī}.

\begin{cousinweb}
\textbf{\ur{لپٹنا}} (\emph{lipaṭnā}) is inherited: CDIAL lists the Hindi cousin \emph{lipaṭnā}, \textquotedblleft{}to stick\textquotedblright{}, under Sanskrit \emph{lipyatē} (CDIAL 11061). The person you cling to takes \textbf{\ur{سے}}: \textbf{\ur{ماں} \ur{سے} \ur{لپٹ} \ur{گئی}}.
\end{cousinweb}

\begin{grammarlens}[title={What you've built}]
One more word for telling what happened.
\end{grammarlens}

\subsection*{Your turn}
\begin{itemize}
\raggedright
  \item \emph{Say it:} \emph{lipaṭnā}
  \item \emph{Say it:} \emph{lipaṭnā}, once more
  \item \emph{Say it:} \emph{bachchī māṅ se lipaṭ gaī.} — The little girl clung to her mother.
  \item \emph{Recall:} say the Urdu for sweet vermicelli, then the Urdu for congratulations, then say \emph{lipaṭnā} again
\end{itemize}

\begin{tcolorbox}[breakable,colback=teal!4,colframe=teal!35!black,title={Before you move on}]
What does \ur{لپٹنا} mean? (\textquotedblleft{}To cling; to hug\textquotedblright{}.) Say it once more.
\end{tcolorbox}
